\section{主干、信息流与完整设计确认}
\label{app:q3}
\begingroup
\renewenvironment{table}[1][]{\begin{readingtable}}{\end{readingtable}}
\renewenvironment{table*}[1][]{\begin{readingtable}}{\end{readingtable}}

\subsection{参数组织与选择}
主干比较固定历史$H$、未来表示$Z$及参照信息流$G_0$。在$G_0$下，未来监督通过共享参数影响动作学习，但Action不读取生成的未来。D的观测/meta、未来和Action共享注意力QKV与输出投影，各有独立FFN。适配后的E分为World和Action两个参数组，每组各有QKV、输出投影、AdaLN和FFN；观测/meta及启用的未来tokens共用World参数。观测/meta令扩散时间$\tau=0$，未来tokens保留各自时间。未来仅由确定性读出提供时，两种主干都不增加闲置生成分支。

内部归一化属于比较变量：E的World AdaLN在$\tau=0$仍有随机初始化的scale/shift偏置，D的观测LayerNorm则从缩放1、偏移0开始。因此，D/E比较的是参数组织与归一化的完整选择。早期E系统的输入和联合生成程序不同，仍作为独立系统参照。

两臂均以种子0、1、2重新训练。同seed的共用输入投影、历史/时间参数和读出头初值相同，主干参数组独立初始化。参照配方为四轮、全局批量16、40,104次更新和641,664个训练窗口（附录B），合计六个端点与8,400场开发评测。D的FFN宽度为$4d$、$d=768$；E在64的正整数倍中选活动主干参数差最小的宽度，同差取较小者，要求差不超过5\%，保持层数、注意力宽度和头数；前端、教师与总参数另列。相同曝光不代表相同FLOPs。

固定参照设备采用原实测NVML标识NVIDIA RPBZZZ6。设备显存为102,641,958,912字节，计算采用FP32并启用TF32。所有候选在该设备和精度下满足首条动作p95耗时不超过200\,ms、峰值显存不超过设备容量80\%。候选减参照的SR配对95\%区间完全高于$+2$个百分点，支持实际重要优势；完全位于$[-2,2]$，支持这一范围内的等效。若候选p95降低至少10\%，且SR区间下界高于$-2$个百分点，可按非劣成本规则选择。每seed成绩及原生任务/安全结果与判定一并报告。三seed不足以判定时，两臂同时增加seed3、4；仍无支配关系则保留参照，不称统计优胜。完整默认参照为T8 RGB与因果估计运动、场景轨迹加目标位姿监督、D主干、R1仅动作部署。使用oracle的设计研究保留其独立输入条件。

% Planned results: em-dash cells are pending measurements, not zeros.
\begin{table*}[t]
\centering\small
\caption{\tablelang{Backbone comparison and directed revisit. SR averages common development scenes and training seeds. Parameter counts cover the active backbone; p95 is time to the first issued action. The revisit fixes the selected information flow.}{主干比较及定向复验。SR汇总共同开发场景及训练种子；参数量仅计活动主干，p95为首条动作前的策略耗时。复验固定在开发集选出的信息流。}}
\label{tab:backbone_confirmation}
\begin{tabular}{@{}llrrrrr@{}}
\toprule
\tablelang{Flow}{信息流}&\tablelang{Backbone}{主干}&\tablelang{Params (M)}{参数量(M)}&SR (\%)&$\Delta$SR (pp)&p95 (ms)&\tablelang{Memory (GiB)}{显存(GiB)}\\
\midrule
$G_0$&D&\resultplaceholder&\resultplaceholder&\tablelang{ref.}{参照}&\resultplaceholder&\resultplaceholder\\
$G_0$&E&\resultplaceholder&\resultplaceholder&\resultplaceholder&\resultplaceholder&\resultplaceholder\\
$G^\star$&D&\resultplaceholder&\resultplaceholder&\tablelang{ref.}{参照}&\resultplaceholder&\resultplaceholder\\
$G^\star$&E&\resultplaceholder&\resultplaceholder&\resultplaceholder&\resultplaceholder&\resultplaceholder\\
\bottomrule
\end{tabular}
\end{table*}


\subsection{共同生成未来接口}
\label{app:generated_future}
各信息流行绑定相同的生成状态：场景位移$F\in\mathbb R^{196\times8\times3}$或稀疏图像轨迹$F\in\mathbb R^{128\times8\times2}$，覆盖50--400\,ms。位移在固定机器人基座系中表示；轨迹保持当前源点身份，图像坐标除以宽、高。源点位置和当前有效性属于观测。未来坐标按有效训练元素逐坐标通道标准化，尺度下限为$10^{-6}$，各行共用统计。每个源点的完整轨迹投影为一个token，加入源点位置和表示类型嵌入，输出投影恢复有序的八个时域。

标准化目标$x_1$与$\epsilon\sim\mathcal N(0,I)$构成$x_\tau=(1-\tau)\epsilon+\tau x_1$，FM目标为$x_1-\epsilon$；每个未来张量共享一个$\tau\sim U[0,1)$。坐标MSE只平均有效元素。无效未来坐标在整个插值中保持其源噪声，不计损失，其掩码不作为模型条件。当前源点的padding在训练与部署中采用相同掩码。独立logit头用有效对应上的BCE预测轨迹可见性，未来可见性GT不进入生成器或Action。目标位姿仍由meta确定性读出，采用附录B的几何损失，不是生成状态$F$；World到Action的边公开未来tokens及其隐状态。

这一归一化高斯接口用于整个信息流比较。较早表征实验采用0.02\,m的米制场景流源和确定性轨迹读出，因此其检查点不能替代新的匹配训练。R1--R6固定辅助权重，Alpha为零。

\subsection{七种信息流程序}
\begin{table*}[t]
\centering\small
\caption{Planned Q3 conditions; no measured result is assigned to these rows.
F denotes a generated future state and A an action state. Arrows specify
information access, independently of whether the future loss is enabled.
R1 has two deployments of the same trained checkpoint. Sequential second
stages use clean ground truth in training and generated first-stage values at test time.}
\label{tab:q3_plan}
\begin{tabularx}{\textwidth}{@{}lcccXX@{}}
\toprule
ID & Future loss & $F\to A$ & $A\to F$ & Cross-stream training input & Inference order \\
\midrule
R0 & No & --- & --- & None & A only \\
R1 & Yes & --- & --- & None & A only \\
R1 & Yes & --- & --- & None & Joint F and A \\
R2 & Yes & Yes & --- & Synchronous noisy states & Joint F and A \\
R3 & Yes & --- & Yes & Synchronous noisy states & Joint F and A \\
R4 & Yes & Yes & Yes & Synchronous noisy states & Joint F and A \\
R5 & Yes & Yes & --- & Clean F for second-stage A & F, then A \\
R6 & Yes & --- & Yes & Clean A for second-stage F & A, then F \\
\bottomrule
\end{tabularx}
\end{table*}

记$C$为观测/meta上下文，$F$为未来tokens，$A$为$8\times7$动作状态。上下文只读取截止当前采集时刻的信息，不读取$F$或$A$；每个生成分支读取$C$与自身。表~\ref{tab:q3_plan}规定额外的跨分支边。这样，未来标签也不能借上下文间接进入Action。禁用边在全部层、全部注意力头上的logit均为$-\infty$，训练与测试使用相同掩码。

R0移除未来损失和tokens，保留动作/上下文接口；R1加入未来监督、关闭双向访问，同一检查点分别仅生成动作或同时生成未被读取的未来。R2开放$F\to A$，R3开放$A\to F$，R4开放双向。训练时F/A使用独立高斯源和独立均匀扩散时间，各自带时间嵌入；部署时两者同时经过$\{0,.2,.4,.6,.8\}$，共五次网络评估。交换的是带噪状态及其隐表示，不是干净未来标签。

R5先在无Action条件下生成F，再固定其解码tokens，通过$F\to A$生成动作。训练分两次：先做无Action的未来FM，再以干净未来GT为条件做动作FM。第二次不穿过第一阶段采样反传，两项损失仍正常更新共享参数。R6顺序相反：先动作FM，再以干净动作做未来FM。部署时第二阶段读取第一阶段预测，每阶段五步Euler，共十次评估。干净教师条件与预测部署条件的差别属于顺序程序本身。R6可在第一阶段后立即下发动作，同时报告首条动作耗时与完整未来生成耗时。

三个seed对应21个训练端点和24个部署条件（包括R1的两种模式），共33,600场开发评测。各行每次重规划执行一条动作，物理限幅相同，同时记录World/Action各自NFE、实际调用、活动参数、曝光、p50/p95动作耗时、显存与物理预测误差。R0/R1比较监督，R1两部署比较未被读取的生成成本，R1--R4比较访问关系，R2/R5及R3/R6比较完整同步/顺序程序。

\begin{table*}[t]
\centering\small
\caption{\tablelang{Information-flow results. Each condition uses three training seeds and the same history/future interface. NFE counts are per modality; joint rows share calls. Forecast error is decoded to mm for 3D displacement or pixels for tracks, fixing representation within a comparison. An absent generated output is marked n/a.}{信息流结果。每个条件使用三个训练种子及共同历史/未来接口。NFE按模态计，同步行共享调用。预测误差在物理解码后计算（三维位移为mm，轨迹为像素），同一比较中表示固定。n/a表示没有对应生成输出。}}
\label{tab:information_flow_results}
\begin{tabular}{@{}llccrrrr@{}}
\toprule
ID&\tablelang{Deployment}{部署}&$\mathrm{NFE}_F$&$\mathrm{NFE}_A$&SR (\%)&\tablelang{Forecast error}{预测误差}&p95 (ms)&\tablelang{Safety (\%)}{安全失败(\%)}\\
\midrule
R0&\tablelang{Action}{动作}&0&5&\resultplaceholder&n/a&\resultplaceholder&\resultplaceholder\\
R1&\tablelang{Action}{动作}&0&5&\resultplaceholder&n/a&\resultplaceholder&\resultplaceholder\\
R1&\tablelang{Joint, unread future}{同步，未来不回读}&5&5&\resultplaceholder&\resultplaceholder&\resultplaceholder&\resultplaceholder\\
R2&\tablelang{Joint $F\to A$}{同步$F\to A$}&5&5&\resultplaceholder&\resultplaceholder&\resultplaceholder&\resultplaceholder\\
R3&\tablelang{Joint $A\to F$}{同步$A\to F$}&5&5&\resultplaceholder&\resultplaceholder&\resultplaceholder&\resultplaceholder\\
R4&\tablelang{Joint, bidirectional}{同步双向}&5&5&\resultplaceholder&\resultplaceholder&\resultplaceholder&\resultplaceholder\\
R5&\tablelang{$F$ then $A$}{先$F$后$A$}&5&5&\resultplaceholder&\resultplaceholder&\resultplaceholder&\resultplaceholder\\
R6&\tablelang{$A$ then $F$}{先$A$后$F$}&5&5&\resultplaceholder&\resultplaceholder&\resultplaceholder&\resultplaceholder\\
\bottomrule
\end{tabular}
\end{table*}


\subsection{干预与交互确认}
\label{app:design_crosschecks}
未来干预固定当前感知、机器人历史、采样数组和动作解码。恒等替换复制原生未来状态；零运动对三维位移置零，对图像轨迹在各时域重复当前点。匹配交换在同任务、物体和可见性层内，用seed17随机选当前目标位姿相差不超过5\,mm/$5^\circ$的其他开发轨迹；源点还须同角色、当前位置相差不超过两像素，双方采用共同有效点/时域。无匹配窗口不纳入并报告覆盖。时间偏移采用$\delta_h\pm50$\,ms，仅保留有标签支持的时间；噪声标准差为开发集预测RMSE的$0,0.5,1,2$倍，各策略共用随机数，保留表示对应的物理坐标。

首次接触前，oracle诊断使用与机器人动作无关的实测外生目标未来。接触后，从同一仿真快照执行固定候选动作序列及外部运动，产生该动作条件下的未来。干预可能改变动作，因此它测量的是对固定候选未来的响应，不能将这一未来当作任意修订动作的oracle。R2/R5建模无动作条件的示范未来，R3/R4/R6则可读取动作；几何有效性不等于动作相容性。任务与安全结果同首条平移和SO(3)动作变化分别报告。

R2/R4在每个Euler步用原训练统计标准化替代干净未来，并以原初噪构造$F_u^I=(1-u)\epsilon_F+uF^I$；重新计算World隐状态后再评估Action，R4保留当前$A\to F$边。恒等控制逐步复制原生中间$F_u$，用于检查原程序精确重建。R5替换第二阶段的干净解码未来并重新嵌入。R0/R1及R3/R6的已下发动作不读取未来，构成阴性对照。标签有效性只定义配对分析域，不作为输入掩码。表~\ref{tab:future_intervention_results}分别报告动作、接触、任务及覆盖。

% Planned results: em-dash cells are pending measurements, not zeros.
\begin{table*}[t]
\centering\small
\caption{\tablelang{Frozen future-interface interventions. Action changes compare first-issued commands under matched noise. Contact errors use actual contact events; non-contact failures stay in SR and are excluded only from conditional contact metrics. Coverage is the fraction of eligible windows/scenes with paired intervention support.}{冻结未来接口干预。动作变化比较相同初噪下的首条下发命令。接触误差使用真实接触事件；无接触失败保留在SR中，只不进入条件接触指标。覆盖率为可进行配对干预的有效窗口/场景比例。}}
\label{tab:future_intervention_results}
\begin{tabular}{@{}lrrrrrr@{}}
\toprule
\tablelang{Intervention}{干预}&$\Delta a$ (mm)&$\Delta R$ ($^\circ$)&\tablelang{Contact (ms)}{接触(ms)}&\tablelang{Pose (mm)}{位姿(mm)}&SR (\%)&\tablelang{Coverage (\%)}{覆盖(\%)}\\
\midrule
\tablelang{Native future}{原生未来}&0&0&\resultplaceholder&\resultplaceholder&\resultplaceholder&\resultplaceholder\\
\tablelang{Identity copy}{恒等复制}&\resultplaceholder&\resultplaceholder&\resultplaceholder&\resultplaceholder&\resultplaceholder&\resultplaceholder\\
\tablelang{Zero motion}{零运动}&\resultplaceholder&\resultplaceholder&\resultplaceholder&\resultplaceholder&\resultplaceholder&\resultplaceholder\\
\tablelang{Matched shuffle}{匹配交换}&\resultplaceholder&\resultplaceholder&\resultplaceholder&\resultplaceholder&\resultplaceholder&\resultplaceholder\\
\tablelang{Time offset $\pm50$ ms}{时间偏移$\pm50$ms}&\resultplaceholder&\resultplaceholder&\resultplaceholder&\resultplaceholder&\resultplaceholder&\resultplaceholder\\
\tablelang{Noise $0.5/1/2\times$ RMSE}{噪声$0.5/1/2\times$RMSE}&\resultplaceholder&\resultplaceholder&\resultplaceholder&\resultplaceholder&\resultplaceholder&\resultplaceholder\\
\tablelang{Pre-contact oracle}{接触前oracle}&\resultplaceholder&\resultplaceholder&\resultplaceholder&\resultplaceholder&\resultplaceholder&\resultplaceholder\\
\tablelang{Candidate-conditioned oracle}{候选动作条件oracle}&\resultplaceholder&\resultplaceholder&\resultplaceholder&\resultplaceholder&\resultplaceholder&\resultplaceholder\\
\bottomrule
\end{tabular}
\end{table*}


确认矩阵包含三项交互：当前RGB/T8 RGB与当前目标位姿/位姿加区间运动交叉，固定选中未来目标；场景流/稀疏轨迹与$G_0$/选中信息流交叉，保留位姿监督；最后在选中信息流下复验D/E。每格使用seed0--2及共同开发场景，完全相同的端点复用。排序反转作为交互报告，并用于选择完整配置。选中的$(H,Z,B,G)$、匹配仅动作消融及最强完整替代方案，在附录A的1,400个新场景上独立确认，不据此重新选择。

\begin{table*}[t]
\centering\small
\caption{\tablelang{Targeted interactions and independent confirmation. History rows fix the future design; future/flow rows retain target-pose supervision. Development contrasts match seeds and scenes. ID confirmation uses 1,400 new scenes after configuration selection.}{定向交互与独立确认。历史部分固定未来设计；未来/信息流部分保持目标位姿监督。开发比较匹配种子与场景；ID确认在配置选定后使用新的1,400场景。}}
\label{tab:design_crosschecks}
\begin{tabularx}{\textwidth}{@{}lXrrr@{}}
\toprule
\tablelang{Comparison}{比较}&\tablelang{Condition}{条件}&SR (\%)&$\Delta$SR (pp)&p95 (ms)\\
\midrule
$H_{\rm RGB}\times H_{\rm motion}$&\tablelang{Current RGB / current pose}{当前RGB/当前位姿}&\resultplaceholder&\tablelang{ref.}{参照}&\resultplaceholder\\
&\tablelang{T8 RGB / current pose}{T8 RGB/当前位姿}&\resultplaceholder&\resultplaceholder&\resultplaceholder\\
&\tablelang{Current RGB / interval motion}{当前RGB/区间运动}&\resultplaceholder&\resultplaceholder&\resultplaceholder\\
&\tablelang{T8 RGB / interval motion}{T8 RGB/区间运动}&\resultplaceholder&\resultplaceholder&\resultplaceholder\\
\midrule
$Z\times G$&\tablelang{Scene flow / $G_0$}{场景流/$G_0$}&\resultplaceholder&\tablelang{ref.}{参照}&\resultplaceholder\\
&\tablelang{Scene flow / $G^\star$}{场景流/$G^\star$}&\resultplaceholder&\resultplaceholder&\resultplaceholder\\
&\tablelang{Sparse tracks / $G_0$}{稀疏轨迹/$G_0$}&\resultplaceholder&\resultplaceholder&\resultplaceholder\\
&\tablelang{Sparse tracks / $G^\star$}{稀疏轨迹/$G^\star$}&\resultplaceholder&\resultplaceholder&\resultplaceholder\\
\midrule
\tablelang{ID confirmation}{ID确认}&\tablelang{Selected complete design}{选定完整设计}&\resultplaceholder&\tablelang{ref.}{参照}&\resultplaceholder\\
&\tablelang{Matched action-only ablation}{匹配的仅动作消融}&\resultplaceholder&\resultplaceholder&\resultplaceholder\\
&\tablelang{Strongest complete alternative}{最强完整替代方案}&\resultplaceholder&\resultplaceholder&\resultplaceholder\\
\bottomrule
\end{tabularx}
\end{table*}


本体/动作旋转消融比较旋转向量与固定轴XYZ Euler角，定义$R=R_z(\gamma)R_y(\beta)R_x(\alpha)$、$\beta\in[-\pi/2,\pi/2]$，其余角取主值。本体为绝对姿态，Action为左乘相对增量$R_{t+1}=\Delta R R_t$；二者同步改变，物理坐标、命令、数据和预算不变。评价使用SR与SO(3)测地误差，避免依赖输出坐标；它与目标位姿读出编码分别比较。

\section{机器人数据预训练、迁移与部署}
\label{app:systems}
\subsection{数据语义与划分}
机器人基座使用DROID示范\cite{droid}，固定选中的历史、未来表示和信息流。训练单位为因果观察窗口及对齐的未来/动作标签；同一episode的全部窗口不跨集合。采集session与场景按组以seed20261018分配到90/5/5的训练/开发/留出集合，组大小取整后的episode数另报。与下游测试场景重叠的组不进入预训练。先划分再筛选，12.5/25/50/100\%嵌套训练子集保持任务/场景分层，只使用本子集归一化统计。episode、有效窗口、独立动作与有效监督分别计数。

源采集时间定义历史与未来。名义50\,ms网格按字段对齐。因果输入只读取已到达前缀：RGB与状态保留不晚于查询的最近已发出值，记录真实年龄及有效性，不用后到样本重建历史。只有离线标签在有效相邻样本间做位置线性插值和最短路径SLERP；命令在记录生效区间内保持。超过两个源采样周期的gap、对象身份或参考系变化不跨越离线插值。每个字段保留源时刻及有效性；重复帧或将名义15\,Hz序列改标20\,Hz不能改变其物理时间。运动学使用实际间隔。

机器人状态与Cartesian命令统一到基座系，先应用标定的法兰到TCP刚性变换，再计算增量。E7包含当前实测TCP到命令TCP的平移差、左乘相对旋转及映射后的夹爪命令；实测开度仍属于状态。训练归一化仅计有效下发动作，重叠窗口不重复计权。越界动作保留拒绝原因，不通过裁剪冒充有效标签。

\subsection{因果输入与离线监督}
DROID没有物体位姿GT。任务指定直接操作的刚体，每个连续片段维持同一实例和对象参考系；目标提示仅从其实际可用时刻起生效。固定前端为SAM~2.1 Hiera-small的首帧目标框与因果mask传播，加FoundationPose RGB-D注册/跟踪及其发布的refiner/scorer权重\cite{sam2,foundationpose}，二者冻结。CAD或独立参考视图必须在episode前获得，不能使用该episode未来重建的几何作为输入先验。每次只处理新获得的前缀，输出当前mask/pose并追加历史，不重写既有输出。

当前深度失效、跟踪丢失和参考系变化进入object有效位；缺失状态使用学习的missing-object嵌入和零有效位，不填入归一化位姿。训练使用相同因果估计及自然发生的缺失，不加入合成输入dropout。缺少物体先验或位姿时，明确的缺失输入路径仍可学习动作/场景监督，位姿损失屏蔽；完整位姿子集另报。目标初始化、注册及逐帧跟踪都计入系统成本。纯RGB外部适配保留独立输入条件，不用oracle深度或位姿补齐。

离线监督读取完整同步片段、DROID发布标定、记录的目标提示及相同独立物体先验。SAM2前后传播mask，双向IoU至少0.5的帧进入FoundationPose逐帧注册及前后跟踪；少于十六个有效深度像素、变换非有限值、双向位姿相差超过20\,mm或$15^\circ$均无效。无物体先验不生成完整pose标签。

离线融合最小化两份pose估计与相邻区间速度变化的Huber目标：平移除以5\,mm，SO(3)角除以$3^\circ$，Huber转折为1；速度与角速度差乘实际$\Delta t$后分别除以5\,mm和$3^\circ$，平滑项权重0.1。gap或初始轨迹速度跳变超过0.2\,m/s、1\,rad/s处去掉平滑边，以保留突变。初值为平移均值和SO(3)测地中点，固定有效域/边，做二十次流形Gauss--Newton迭代。教师不重写已发出输入。未来目标使用离线轨迹内部差$\tilde p_{t+h}-\tilde p_t$及$\Log(\tilde R_{t+h}\tilde R_t^\top)^\vee$，不把离线未来减去当前因果估计，避免混入当前位姿纠偏。每个时域要求当前参考、未来位姿、身份/参考系和旋转可观测性均有效；缺失监督只屏蔽对应损失。完整SE(3)池按普通SO(3)损失使用可观测刚体位姿；部分旋转或对称等价类不冒充完整标签。

质量核验采用40个片段，14开发、26留出，每片五个锚点。七类覆盖静止、平滑平移、旋转、接触/滑移、遮挡/恢复、弱外观，以及对称或刚性关节部件，各层数量另报。留出片段不调估计器。同前缀截断/替换未来检查已发出位姿、有效位、目标身份和参考几何在重复运行数值容差内不变。有足够约束的片段使用独立人工关键点及标定多视角几何作参考，其余标记物理精度未验证。

因果与离线路径分别报告平移/旋转median、p95，逐时域运动误差、身份失败、错误接受、保留率、恢复时间及计算/标注成本。接触误差预算为$e_p+2r\sin(e_R/2)+v e_t$，其中$r$为接触点力臂，$v$为相对接近速度；在开发阶段根据任务间隙与计时预算确定容差，再用于留出核验。该式是工程误差界，不是联合概率保证。质量与覆盖同时报告。

场景轨迹用冻结DPFlow估计相邻源帧双向flow，双线性复合将固定当前源点传播至八个物理时刻\cite{dpflow}，前后向一致性采用附录C的像素判据；链路中全部坐标及反向检查有效才记对应有效。MoGe-2提供米制伪深度\cite{moge2}，按缩放后的内参与发布外参反投影得到三维位移。伪深度不当作测量GT；可见性与对应有效性分别保存，深度无效则屏蔽三维监督，不制造零运动。各比较共用源点、时间映射和教师版本。

\begin{table*}[t]
\centering\small
\caption{\tablelang{Robot-data accounting. Splits are grouped before filtering. Coverage uses all candidate windows as its denominator; future validity is counted separately at each horizon. Causal-input accuracy and offline-label accuracy use independent references where available.}{机器人数据统计。先按组划分，再筛选窗口。覆盖率以全部候选窗口为分母；未来有效性逐时域统计。因果输入和离线标签的精度分别使用可获得的独立参考。}}
\label{tab:pretraining_data}
\begin{tabular}{@{}lrrrrr@{}}
\toprule
\tablelang{Partition}{划分}&\tablelang{Episodes}{Episode数}&\tablelang{Windows}{窗口数}&\tablelang{Valid actions}{有效动作}&\tablelang{Pose coverage (\%)}{位姿覆盖(\%)}&\tablelang{Track coverage (\%)}{轨迹覆盖(\%)}\\
\midrule
\tablelang{Train (90\% groups)}{训练(90\%组)}&\resultplaceholder&\resultplaceholder&\resultplaceholder&\resultplaceholder&\resultplaceholder\\
\tablelang{Development (5\%)}{开发(5\%)}&\resultplaceholder&\resultplaceholder&\resultplaceholder&\resultplaceholder&\resultplaceholder\\
\tablelang{Held out (5\%)}{留出(5\%)}&\resultplaceholder&\resultplaceholder&\resultplaceholder&\resultplaceholder&\resultplaceholder\\
\midrule
\multicolumn{2}{l}{\tablelang{Held-out pose quality}{留出位姿质量}}&
\multicolumn{2}{c}{\tablelang{Translation med./p95 (mm)}{平移median/p95(mm)}}&
\multicolumn{2}{c}{\tablelang{Rotation med./p95 ($^\circ$)}{旋转median/p95($^\circ$)}}\\
\multicolumn{2}{l}{\tablelang{Causal input}{因果输入}}&\multicolumn{2}{c}{\resultplaceholder}&\multicolumn{2}{c}{\resultplaceholder}\\
\multicolumn{2}{l}{\tablelang{Offline teacher}{离线教师}}&\multicolumn{2}{c}{\resultplaceholder}&\multicolumn{2}{c}{\resultplaceholder}\\
\bottomrule
\end{tabular}
\end{table*}


\subsection{训练配方与规模}
Benchmark参照为$\eta_0=2\times10^{-4}$、$B_0=16$、$T_0=641,664$个有效消费窗口，优化器、裁剪和完整cosine日程见附录B。共同的24候选搜索取学习率倍数$\{0.5,1,2\}$、batch倍数$\{1,4\}$、曝光倍数$\{1,2\}$及既有辅助loss倍数$\{0.5,1\}$；预热消费16,000个样本。每个候选使用自身完整预算的cosine日程，低学习率尾段延长不替代短预算对照。全部24候选先消费各自预算25\%，按开发SR保留六个接续完整日程，前两个各作三seed确认；初始化对照使用相同搜索额度。附录H的实测配方研究保留种子与逐任务波动。

机器人预训练的一次完整pass消费全部有效训练窗口。参照$T_{\rm robot}$为四pass，batch256，其余AdamW设置相同；前2.5\%曝光预热，再cosine衰减到峰值的0.1。在参照规模上用相同24候选网格搜索学习率、batch、曝光及既有loss权重，所选设置定义完整数据预算。这里计有效窗口，不计原始视频帧数。

规模约为参照活动主干的$0.5,1,2,4$倍：固定宽768、8头，采用2/4/8/16层，精确参数量包含投影和读出。表示、分辨率及token组织不变。四档均用全数据；1倍与2倍额外使用三个较小嵌套子集，共十条主曲线。固定曝光时，小子集重复更多次。每档的六次局部试验固定为学习率倍数$\{0.5,1,2\}$与World loss倍数$\{0.5,1\}$的交叉，batch与曝光沿用选中参照。

计算量比较在1倍和2倍模型上使用参照完整计算预算的$0.5,1,2$倍，优先匹配实测训练FLOPs，否则报告指定设备的GPU小时。不共用日程的中间端点不充作匹配预算结果。各完整训练的25/50/75/100\%曝光处评测迁移。源域loss描述学习，开发迁移SR与部署约束选择规模及同一基座端点。每个候选在两域各用一个参照后训练曝光、相同训练数据及seed0--2；满足部署约束后，按仿真/真机各域任务宏平均SR等权排序。真机开发每任务/速度/延迟格每seed使用五个独立物体/初态trial，不与最终测试共享。参照规模及其主要仅动作预训练消融各用三个独立预训练seed；更大规模的一seed曲线明确标注，从一个基座派生三个后训练seed只衡量后训练波动。

% Planned results: em-dash cells are pending measurements, not zeros.
\begin{table*}[t]
\centering\small
\caption{\tablelang{Ten principal model/data curves. Depth changes at fixed width, heads, and tokenization. Data fractions are nested independent-episode subsets; consumed training exposure is common. Each transfer score uses independently initialized downstream branches from that row's base.}{模型/数据的十条主曲线。仅改变深度，固定宽度、注意力头和token组织。数据比例对应嵌套的独立episode子集，累计训练曝光相同。每行的迁移成绩均从该行基座独立建立后训练分支。}}
\label{tab:scale_results}
\begin{tabular}{@{}lccrrrr@{}}
\toprule
\tablelang{Scale}{规模}&\tablelang{Depth}{层数}&\tablelang{Data (\%)}{数据(\%)}&\tablelang{Params (M)}{参数(M)}&\tablelang{Sim SR (\%)}{仿真SR(\%)}&\tablelang{Real SR (\%)}{真机SR(\%)}&\tablelang{GPU h}{GPU小时}\\
\midrule
$0.5\times$&2&100&\resultplaceholder&\resultplaceholder&\resultplaceholder&\resultplaceholder\\
$1\times$&4&12.5&\resultplaceholder&\resultplaceholder&\resultplaceholder&\resultplaceholder\\
$1\times$&4&25&\resultplaceholder&\resultplaceholder&\resultplaceholder&\resultplaceholder\\
$1\times$&4&50&\resultplaceholder&\resultplaceholder&\resultplaceholder&\resultplaceholder\\
$1\times$&4&100&\resultplaceholder&\resultplaceholder&\resultplaceholder&\resultplaceholder\\
$2\times$&8&12.5&\resultplaceholder&\resultplaceholder&\resultplaceholder&\resultplaceholder\\
$2\times$&8&25&\resultplaceholder&\resultplaceholder&\resultplaceholder&\resultplaceholder\\
$2\times$&8&50&\resultplaceholder&\resultplaceholder&\resultplaceholder&\resultplaceholder\\
$2\times$&8&100&\resultplaceholder&\resultplaceholder&\resultplaceholder&\resultplaceholder\\
$4\times$&16&100&\resultplaceholder&\resultplaceholder&\resultplaceholder&\resultplaceholder\\
\bottomrule
\end{tabular}
\end{table*}


\subsection{独立后训练与消融}
仿真与真机分别直接从同一机器人基座初始化，各自建立优化器、日程、训练归一化与任务读出；仿真权重不进入真机分支。不兼容的动作/状态投影按相同seed重新初始化，其余兼容参数直接载入。有无机器人预训练均保留相同通用视觉初始化与下游数据，分别报告共同配方和同等24候选调参后的表现。示范比例为嵌套的10/25/50/100\%，各档保留任务覆盖。仿真参照采用Benchmark配方；真机为四pass有效示范。两域各按相同相对网格选择配方，并使用seed0--2。下游预算相同不表示总计算相同，预训练成本单列。

纯动作预训练使用R0，只迁移共用参数，新增未来参数采用相同初值。另一项消融保留全部通路，仅移除显式World loss；当Action读取未来时，动作梯度或顺序教师强制仍可使用该通路，因此不能称为无未来信息对照。其余比较包括：World监督仅在后训练加入/两阶段均加入；选中伪标签筛选与有效标签数量相等的未筛选质量对照；冻结/训练兼容视觉参数。预训练后定向复核最关键的两项历史/未来干预。仅在纳入额外真实来源时，比较相同独立数据量与曝光下的单源/多源。内部仿真、外部基准与真机结果分别报告。

\begin{table*}[t]
\centering\small
\caption{\tablelang{Transfer and pretraining controls at the selected scale. All rows share downstream demonstrations and matched search allocation. Simulation, native external suites, and physical tasks retain distinct score denominators. The four LIBERO entries are Spatial/Object/Goal/Long.}{选定规模的迁移与预训练对照。各行共享下游示范及相同调参预算。内部仿真、原生外部基准和真机分别保留评分分母。LIBERO四项依次为Spatial/Object/Goal/Long。}}
\label{tab:transfer_results}
\begin{tabularx}{\textwidth}{@{}Xrrrr@{}}
\toprule
\tablelang{Initialization / objective}{初始化/目标}&\tablelang{Internal SR}{内部SR}&DOMINO&LIBERO&\tablelang{Real SR}{真机SR}\\
\midrule
\tablelang{No robot pretraining}{无机器人预训练}&\resultplaceholder&\resultplaceholder&\resultplaceholder&\resultplaceholder\\
\tablelang{Action-only robot pretraining (R0)}{仅动作机器人预训练(R0)}&\resultplaceholder&\resultplaceholder&\resultplaceholder&\resultplaceholder\\
\tablelang{Complete robot pretraining}{完整机器人预训练}&\resultplaceholder&\resultplaceholder&\resultplaceholder&\resultplaceholder\\
\tablelang{Without explicit world loss; paths retained}{去显式World损失，保留通路}&\resultplaceholder&\resultplaceholder&\resultplaceholder&\resultplaceholder\\
\tablelang{Matched label-count quality control}{标签数匹配的质量对照}&\resultplaceholder&\resultplaceholder&\resultplaceholder&\resultplaceholder\\
\tablelang{Frozen visual encoder}{冻结视觉编码器}&\resultplaceholder&\resultplaceholder&\resultplaceholder&\resultplaceholder\\
\midrule
\tablelang{Downstream data 10\% / 25\% / 50\% / 100\%}{下游数据10\%/25\%/50\%/100\%}&\resultplaceholder&\resultplaceholder&\resultplaceholder&\resultplaceholder\\
\bottomrule
\end{tabularx}
\end{table*}


\subsection{外部基准与真机任务}
单臂DOMINO\cite{domino}采用原生动态拦截/跟踪任务、官方训练/开发划分及成功评测器；每任务每训练seed使用50个留出场景，原生任务数随结果报告，任务等权平均，不替换成内部十四任务的阈值。LIBERO\cite{libero}采用Spatial、Object、Goal、Long四suite，每suite十任务，每任务50个评测场景，三个训练seed；示范与初态划分遵循发布协议，四suite成绩分别列出。

参照硬件契约为七自由度力矩控制机械臂、平行夹爪、标定的固定RGB-D相机及带时间戳的Cartesian TCP接口；这些是平台要求，型号和实测标定作为部署元数据。几何依据独立测量的物体尺度及容器/开口间隙，各方法在试验前固定同一装置尺寸。每trial自目标开始运动起至多15\,s。抓取要求提升至初始支撑面上方至少50\,mm并保持1\,s；放置要求完全释放且全部物体几何在容器内保持1\,s；穿越要求全部物体越过开口平面、无障碍碰撞，再保持1\,s；再次拦截要求外部撞击后完成相同提升/保持，撞击前机器人接触即失败。

真机包含移动目标抓取、移动容器放置、限时开口穿越、外部碰撞后再次拦截四类。每类使用三档速度$(0.5,1,1.5)v_{\rm ref}$和三档新增命令延迟$(0,1,2)\Delta_c$。以装置声明速度上限为界，做十步二分搜索，取连续十次无装置碰撞、越界或固定结构失稳的独立运动标定循环所能通过的最高速度；$v_{\rm ref}$为其三分之二。$\Delta_c$为实测控制周期。延迟在推理之后增加，与总动作年龄一并报告。每格十个物体/初态试验，每seed为$4\times3\times3\times10=360$次，三个seed合计每方法1,080次，同条件逐方法配对。机械臂、夹爪、相机、控制周期、标定、物体尺寸和绝对速度均作为部署元数据填写，不视为已知硬件事实。

最终测试的物体和背景布局与下游训练/开发分离。固定循环方法顺序在相邻trial间轮换，以分散漂移；每个方法结束后重置。人工接触、安全停止或手动完成任务计失败。首次有效观察前的基础设施故障对整个配对block重试并单列。成功要求相应物理条件维持规定时长；碰撞、掉落、越界、介入、超时分别记录。接触时间误差相对任务允许窗口，位姿误差在首次接触时计算。没有接触的失败保留在SR分母，接触误差记缺失，不填零。

\subsection{延迟与成本}
固定延迟扫描在原生命令队列上叠加0/50/100/150/200\,ms，主比较取新增150\,ms。各方法使用相同观察截止和队列，仿真物理在规定延迟内继续演化。它与包含采集、预处理、推理、排队、执行的实际部署分别报告。

\begin{table*}[t]
\centering\small
\caption{\tablelang{Delay sensitivity and physical deployment cost. The controlled delay is added to the native command queue; synchronous microbenchmark latency is not substituted for action age. SE(3)-WAM and its matched R0 control use the same estimator and device.}{延迟敏感性与物理部署成本。受控延迟叠加在原生命令队列上；同步计时不能替代动作年龄。SE(3)-WAM与匹配R0对照使用相同估计器和设备。}}
\label{tab:delay_results}
\begin{tabular}{@{}lrrrrrrr@{}}
\toprule
&\multicolumn{5}{c}{\tablelang{SR (\%) at added delay (ms)}{不同新增延迟(ms)下SR(\%)}}&
\tablelang{Age p95}{年龄p95}&\tablelang{Miss (\%)}{超期(\%)}\\
\cmidrule(lr){2-6}
\tablelang{Method}{方法}&0&50&100&150&200&(ms)&\\
\midrule
\tablelang{Matched R0}{匹配R0}&\resultplaceholder&\resultplaceholder&\resultplaceholder&\resultplaceholder&\resultplaceholder&\resultplaceholder&\resultplaceholder\\
SE(3)-WAM&\resultplaceholder&\resultplaceholder&\resultplaceholder&\resultplaceholder&\resultplaceholder&\resultplaceholder&\resultplaceholder\\
\tablelang{Strongest system reference}{最强系统参照}&\resultplaceholder&\resultplaceholder&\resultplaceholder&\resultplaceholder&\resultplaceholder&\resultplaceholder&\resultplaceholder\\
\bottomrule
\end{tabular}
\end{table*}


每个设备/精度组合先预热三十次，再在相同三条T8序列上计时300次，调用边界同步CUDA。三轮计时轮换方法顺序。首条动作耗时包含内部感知与动作生成，完整预测另报；外部图像I/O和渲染不计入该微基准。部署记录采集、可用、预处理、推理起止、入队、下发及应用时刻，动作年龄为应用减源采集时间。报告年龄p50/p95/max、超期比例、峰值allocated与总常驻显存，以及设备、精度、batch和采样调用。选择限值为p95动作耗时不超过200\,ms、显存不超过设备容量80\%。

原未来监督实验的策略调用均值177.071--184.212\,ms，p95为178.726--185.133\,ms，峰值allocation为1.571--1.617\,GiB。该原协议使用六次预热、三序列共三十次测量，每次reset后七次真实调用priming，再计第八次；FP32/TF32、$B=S=1$、五步Euler、CPU affinity0--2。同设备NVML名为NVIDIA RPBZZZ6，显存102,641,958,912字节；加载、外部oracle、渲染、I/O和priming不计时。这些描述性数据保留原采样口径。

\subsection{初始化控制}
DreamVLA对照仅去掉LIBERO策略检查点，保留MAE/CLIP、教师目标、损失和独立二十轮配方（48,220更新、batch512）。$\pi_{0.5}$将base改为LIBERO expert，保留30,000更新、batch256；两者使用原十四任务划分。

第三项用仅动作D（四层、宽768）、可训练CNN、T8 RGB/机器人历史，去掉目标位姿、深度/相机几何、未来流及目标区域/未来损失。从头与LIBERO-90初始化版本共用11,200条示范拟合的归一化。LIBERO预训练为4,050训练、450留出示范，三十轮、8,520更新、batch256；下游均重置优化器/日程，四轮、40,104更新、batch16。这是简化策略的视觉运动迁移，独立于正式机器人数据基座和语言条件LIBERO评测。

\begin{table*}[t]
\centering\small
\caption{\tablelang{Auxiliary initialization controls on the original fourteen-task development partitions. Training exposure is fixed within each pair; these are separate from the robot-data-base study.}{原十四任务开发划分上的辅助初始化对照。每对固定训练曝光，与正式机器人数据基座研究分别报告。}}
\label{tab:initialization_controls}
\begin{tabularx}{\textwidth}{@{}lXrr@{}}
\toprule
\tablelang{Policy}{策略}&\tablelang{Initialization}{初始化}&\tablelang{140-scene SR (\%)}{140场SR(\%)}&\tablelang{1,260-scene SR (\%)}{1,260场SR(\%)}\\
\midrule
DreamVLA&\tablelang{LIBERO policy; MAE/CLIP}{LIBERO策略；MAE/CLIP}&\resultplaceholder&\resultplaceholder\\
DreamVLA&\tablelang{MAE/CLIP only}{仅MAE/CLIP}&\resultplaceholder&\resultplaceholder\\
$\pi_{0.5}$&\tablelang{Base}{基础初始化}&\resultplaceholder&\resultplaceholder\\
$\pi_{0.5}$&\tablelang{LIBERO expert}{LIBERO专家}&\resultplaceholder&\resultplaceholder\\
\tablelang{Action-only D}{仅动作D}&\tablelang{Fresh}{从头训练}&\resultplaceholder&\resultplaceholder\\
\tablelang{Action-only D}{仅动作D}&\tablelang{LIBERO-90 visual-motor}{LIBERO-90视觉运动}&\resultplaceholder&\resultplaceholder\\
\bottomrule
\end{tabularx}
\end{table*}

\endgroup
